% Adaptación de FORMATO.tex para AA2; plantilla original preservada.
\documentclass[12pt,a4paper]{report}

% ------------------------------------------
%   PAQUETES BÁSICOS
% ------------------------------------------
\usepackage[spanish]{babel}
\usepackage[utf8]{inputenc}
\usepackage[T1]{fontenc}
\usepackage{graphicx}
\usepackage{amsmath, amssymb}
\usepackage{booktabs}
\usepackage{geometry}
  \geometry{left=3cm, right=2.5cm, top=3cm, bottom=2.5cm}
\usepackage{setspace}
  \onehalfspacing
\usepackage{titlesec} 
  \titleformat{\chapter}[block]{\bfseries\large}{\thechapter.}{1em}{}
  \titleformat{\section}[block]{\bfseries}{\thesection.}{1em}{}
\usepackage{tocloft}
  \setlength{\cftbeforechapskip}{6pt}
  \renewcommand{\cftchapfont}{\bfseries}
\usepackage{longtable}
\usepackage{array}
\usepackage{verbatim}

\usepackage{cite}
\usepackage{hyperref}
% Soporte mínimo para tablas, rutas y bloques técnicos.
\usepackage{calc}
\usepackage{xurl}
\usepackage{fvextra}
\providecommand{\tightlist}{\setlength{\itemsep}{0pt}\setlength{\parskip}{0pt}}
\setlength{\emergencystretch}{3em}

\addto\captionsspanish{
    \renewcommand{\contentsname}{Índice General}
    \renewcommand{\tablename}{Tabla}
}

% ------------------------------------------
%   FONDO DE PÁGINA (SOLO CARÁTULA)
% ------------------------------------------
\usepackage{eso-pic}
\newcommand\BackgroundPic{%
  \put(0,0){%
    \parbox[b][\paperheight]{\paperwidth}{%
      \vfill
      \centering
      \IfFileExists{fondo.jpeg}{\includegraphics[width=\paperwidth,height=\paperheight]{fondo.jpeg}}{\parbox{0.8\paperwidth}{\centering\small [PENDIENTE: fondo.jpeg original]}}%
      \vfill
}}}

% ------------------------------------------
%   ENCABEZADO CON LOGO (OTRAS PÁGINAS)
% ------------------------------------------
\usepackage{fancyhdr}
\pagestyle{fancy}
\fancyhf{}
\setlength{\headheight}{40pt}
\fancyhead[L]{\IfFileExists{logo_unl.jpeg}{\includegraphics[height=1.2cm]{logo_unl.jpeg}}{\fbox{\scriptsize PENDIENTE: logo\_unl.jpeg}}}
\fancyfoot[C]{\thepage}
\renewcommand{\headrulewidth}{0pt}

\fancypagestyle{plain}{
    \fancyhf{}
    \fancyhead[L]{\IfFileExists{logo_unl.jpeg}{\includegraphics[height=1.2cm]{logo_unl.jpeg}}{\fbox{\scriptsize PENDIENTE: logo\_unl.jpeg}}}
    \fancyfoot[C]{\thepage}
    \renewcommand{\headrulewidth}{0pt}
}

% ------------------------------------------
%   METADATOS DEL DOCUMENTO
% ------------------------------------------
\newcommand{\TituloTrabajo}{AA2 -- Dossier de requisitos y matriz de trazabilidad}
\newcommand{\Autor}{[PENDIENTE: estudiante]}
\newcommand{\Director}{Edison Leonardo Coronel Romero (referencia APE2)}
\newcommand{\Carrera}{Carrera de Computación}
\newcommand{\Facultad}{Facultad de la Energía, las Industrias y los Recursos Naturales No Renovables}
\newcommand{\Universidad}{Universidad Nacional de Loja}
\newcommand{\Ciudad}{Loja, Ecuador}
\newcommand{\Fecha}{Semana 2 -- [PENDIENTE: fecha real]}

\begin{document}

\AddToShipoutPictureBG*{\BackgroundPic}

% ------------------------------------------
%   PORTADA
% ------------------------------------------
\begin{titlepage}
  \centering
  \vspace*{2cm}
  {\bfseries\Large \Universidad}\\[0.5cm]
  {\bfseries\large \Facultad}\\[0.5cm]
  {\bfseries\large \Carrera}\\[2cm]
  {\bfseries\Large \TituloTrabajo}\\[2.5cm]

  \begin{flushright}
    \setlength{\fboxsep}{8pt}
    \fbox{%
      \begin{minipage}{0.4\textwidth}\small
        Desarrollo Basado en Plataformas (referencia APE2).
      \end{minipage}%
    }
  \end{flushright}
  \vspace{1cm}

  {\bfseries \large AUTOR:}\\[0.5em]
  {\large \Autor}\\[2em]
  {\bfseries \large TUTOR:}\\[0.5em]
  {\large \Director}\\[2em]
  {\bfseries \large \Ciudad}\\[1em]
  {\bfseries \large \Fecha}
\end{titlepage}

% ------------------------------------------
%   DECLARACIÓN DE USO DE IA GENERATIVA
% ------------------------------------------
\newpage
\chapter*{Declaración de Delegación a IA Generativa}
\addcontentsline{toc}{chapter}{Declaración de uso de IA Generativa}
Se utilizó ChatGPT/Codex para revisar documentación y fuentes técnicas y del dominio, proponer depuración de historias y requisitos, construir la matriz y generar los documentos. No se realizaron entrevistas, validación docente, pruebas de negocio, solicitudes al backend ni publicación.

El estudiante debe revisar el contenido, identificar su aporte real y asumir responsabilidad por su entrega. No se declara supervisión humana ya completada ni se reutilizan autores o herramientas del informe previo contenido en la plantilla.

\noindent\textbf{Declaración enviada por: [PENDIENTE: estudiante real].}
\newpage
\tableofcontents
\clearpage
\chapter*{Ficha de la actividad}
\addcontentsline{toc}{chapter}{Ficha de la actividad}
\textbf{Proyecto:} SIGEAC / SEC-UNL.\par
\textbf{Actividad:} AA2, Semana 2, Actividad Autónoma.\par
\textbf{Modalidad:} trabajo individual articulado con el proyecto grupal.\par
\textbf{Duración prevista:} 4,5 horas; no es tiempo de ejecución medido.\par
\textbf{Asignatura, docente y ciclo:} referencias de continuidad con APE2, pendientes de confirmación si esta guía corresponde a otra asignatura.\par
\textbf{Base técnica:} commit \texttt{faf868ea5a8a85285290c06d3baf8fa49889edac} y especificación académica local APE2.\par
\textbf{Propósito:} justificar el origen de cada requisito y cómo se verificará, vinculando actor, necesidad, comportamiento, criterio y fuente.

El dossier incluye la matriz, referencias y nota individual como anexos. Las fuentes primarias consultadas no certifican vigencia ni competencias del piloto. La revisión individual y validación de dominio siguen pendientes.
\chapter{Revisión de fuentes y fundamentos
técnicos}\label{dossier-revisiuxf3n-de-fuentes-y-fundamentos-tuxe9cnicos}

\section{Método y criterios de
depuración}\label{dossier-muxe9todo-y-criterios-de-depuraciuxf3n}

Se contrastaron las tres historias y los diez FR/diez NFR de APE2 con
las necesidades del proyecto, documentos primarios y fuentes
ejecutables. Se separaron comportamientos compuestos, se explicitaron
límites de las fuentes y se retiraron de la selección activa compromisos
cuya pertinencia al recorte no estaba justificada suficientemente.
``Retirar de la selección'' no elimina el requisito global ni significa
aprobación de un cambio de alcance.

Cada elemento se evalúa por claridad, necesidad, consistencia,
verificabilidad, ausencia de redundancia y compatibilidad con su fuente.
La revisión diferencia decisiones del producto, fundamento técnico,
contexto institucional e implementación observada. Ninguna fuente se
utiliza para atribuir una entrevista, designación o aprobación
inexistente.

\section{Fuentes consultadas y alcance de la
evidencia}\label{dossier-fuentes-consultadas-y-alcance-de-la-evidencia}

El registro completo F-01--F-14 está en \url{FUENTES_REFERENCIAS.md},
con autor, localizador, relación y límite. Los principales insumos son
APE2, especificación y reglas del equipo, los PDFs oficiales CACES/UNL,
RFC 9110 y WCAG 2.2.

El
\href{https://www.caces.gob.ec/wp-content/uploads/2024/04/Modelo-generico-para-la-evaluacio\%CC\%81n-del-entorno-de-aprendizaje-de-carreras-de-grado.pdf}{Modelo
Genérico CACES 2024}, secciones II y V, respalda contextualizar la
evaluación mediante modelo, fuentes y períodos. El
\href{https://www.unl.edu.ec/sites/default/files/gaceta/resolucion/2022-06/Resoluci\%C3\%B3n\%20Nro.\%20037-R-UNL-2022-\%20Instructivo\%20para\%20el\%20Aseguramiento\%20de\%20la\%20Calidad\%20de\%20la\%20UNL-signed.pdf}{Instructivo
UNL 037-R-UNL-2022}, artículos 9--12, aporta un correlato institucional
para funciones de evaluación.

Estas consultas permiten avanzar en procedencia documental respecto a
APE2. No corroboran las configuraciones sintéticas ni la vigencia actual
de esas disposiciones para el piloto. F-10 mantiene expresamente
\path|institutionalValidity: false|; no se altera ese registro. El
vínculo entre órganos institucionales y roles del software continúa
pendiente de designaciones y validación.

\section{Lectura de RFC 9110 aplicada al
proyecto}\label{dossier-lectura-de-rfc-9110-aplicada-al-proyecto}

Síntesis de la \href{https://www.rfc-editor.org/rfc/rfc9110.html}{RFC
9110}, organizada por los seis conceptos solicitados:

{\def\LTcaptype{none} % do not increment counter
\begin{longtable}[]{@{}
  >{\raggedright\arraybackslash}p{(\linewidth - 2\tabcolsep) * \real{0.2500}}
  >{\raggedright\arraybackslash}p{(\linewidth - 2\tabcolsep) * \real{0.7500}}@{}}
\toprule\noalign{}
\begin{minipage}[b]{\linewidth}\raggedright
Concepto / localizador
\end{minipage} & \begin{minipage}[b]{\linewidth}\raggedright
Síntesis técnica
\end{minipage} \\
\midrule\noalign{}
\endhead
\bottomrule\noalign{}
\endlastfoot
Métodos, §§9.2--9.3 & GET recupera; POST procesa; PUT expresa reemplazo;
DELETE solicita retirar una asociación. Seguridad e idempotencia son
propiedades distintas. \\
Representación, §§8.1--8.3 & El recurso se comunica mediante
representaciones y metadatos; Content-Type identifica el tipo. \\
Estados, §15 & El estado distingue éxito, rechazo y error; 201 indica
creación y 412 una precondición incumplida. \\
Cabeceras, §§5, 8.8.3, 10.2.2, 12.5.1 & ETag identifica una
representación; Location referencia el recurso; Accept expresa
preferencias. \\
Contenido, §6.4 & El significado depende de método y estado; 204 no
incluye contenido. \\
Seguridad, §§11 y 17 & Considerar autenticación, divulgación e
intermediarios; HTTP por sí solo no resuelve la autorización de
negocio. \\
\end{longtable}
}

La definición de PATCH se consulta aparte en F-14. Estas fuentes no
fijan los 30 minutos de sesión ni el p95 de 2 s; ambos son decisiones
del proyecto F-03.

\section{Decisiones y límites observados en la
API}\label{dossier-decisiones-y-luxedmites-observados-en-la-api}

El código C-01 establece sesión con cookie y entrega CSRF; no se deduce
soporte JWT/Bearer del ejemplo académico. C-03 separa alta y
modificación de carrera. C-05 crea procesos en preparación y no permite
que el cliente envíe un estado de activación en su DTO de alta.

En el contrato actual, editar una carrera utiliza PATCH con CareerInput
completo; no debe enviarse un subconjunto arbitrario por creer que el
método lo admite. La revisión empieza en 1 para carrera y en 0 para
proceso. Esas diferencias son de los DTO/controladores, no reglas
universales de HTTP.

C-06 acepta un ETag fuerte numérico y produce 400/428 ante formatos o
precondición ausente según su contrato. La gramática general de If-Match
es más amplia; no se califica una lista o \path|*| como inválida para
todo HTTP por la restricción local. Referencia de contraste: F-09
§13.1.1.

\textbf{Observación HTTP-AA-01, pendiente:} RFC 9110 §15.5.2 exige un
desafío WWW-Authenticate en 401. C-07/C-08 no lo establecen
explícitamente y no se localizó la cabecera en fuentes principales.
Falta observar respuestas reales y decidir el desafío aplicable; no se
incorpora un Bearer ficticio ni se afirma conformidad integral de
autenticación.

La aceptación de los contratos preliminares CP-01--09 de APE2 queda
condicionada a ese contraste. Los ejemplos son reconstruidos y no se
ejecutaron solicitudes.

\chapter{Dossier depurado del
incremento}\label{dossier-dossier-depurado-del-incremento}

\section{Síntesis del problema y
alcance}\label{dossier-suxedntesis-del-problema-y-alcance}

El proyecto organiza información y decisiones de autoevaluación interna.
Para iniciar un proceso necesita identificar carrera, unidad, período,
versión de modelo y actores autorizados, conservando el historial sin
atribuir acreditación. No se aportó un diagnóstico empírico que mida
pérdidas de documentos o tiempos institucionales; tales cifras no se
inventan.

Se conserva el recorte propuesto en APE2: acceso, catálogo de
unidades/carreras y preparación/consulta de un proceso con modelo
publicado. No incluye autorización/activación institucional, evidencias,
revisión, valoración, informes, mejora ni archivo. Su aceptación docente
sigue pendiente en CL-APE-01.

Las decisiones de tecnología se mantienen fuera de las necesidades: las
fuentes muestran Angular, Spring Boot y persistencia, pero el actor
necesita una operación identificable y autorizada, no una clase o un
framework concreto.

\section{Actores y necesidades
justificadas}\label{dossier-actores-y-necesidades-justificadas}

{\def\LTcaptype{none} % do not increment counter
\begin{longtable}[]{@{}
  >{\raggedright\arraybackslash}p{(\linewidth - 6\tabcolsep) * \real{0.2500}}
  >{\raggedright\arraybackslash}p{(\linewidth - 6\tabcolsep) * \real{0.2500}}
  >{\raggedright\arraybackslash}p{(\linewidth - 6\tabcolsep) * \real{0.2500}}
  >{\raggedright\arraybackslash}p{(\linewidth - 6\tabcolsep) * \real{0.2500}}@{}}
\toprule\noalign{}
\begin{minipage}[b]{\linewidth}\raggedright
ID
\end{minipage} & \begin{minipage}[b]{\linewidth}\raggedright
Actor documentado
\end{minipage} & \begin{minipage}[b]{\linewidth}\raggedright
Necesidad
\end{minipage} & \begin{minipage}[b]{\linewidth}\raggedright
Fuente y límite
\end{minipage} \\
\midrule\noalign{}
\endhead
\bottomrule\noalign{}
\endlastfoot
A-AA-01 & Usuario habilitado & N-AA-01: conocer su acceso vigente y
terminarlo cuando concluya su trabajo. & F-03/F-05; identidad sintética
disponible, persona real no identificada. \\
A-AA-02 & Administrador funcional autorizado & N-AA-02: mantener
unidades/carreras identificables y con historial. & F-03/F-06;
competencia institucional exacta por confirmar. \\
A-AA-03 & Coordinador con capacidad de gestión en la carrera & N-AA-03:
preparar una evaluación con referencia exacta de carrera, período y
modelo. & F-03/F-04; contexto F-07/F-08. No se equipara a una comisión
por inferencia. \\
A-AA-04 & Consultor o lector autorizado & N-AA-04: consultar únicamente
recursos visibles y comprender su estado. & F-03/F-05; los permisos de
catálogo y proceso no son intercambiables. \\
A-AA-02/03 & Actores que modifican recursos & N-AA-05: evitar pérdida de
cambios y poder explicar qué ocurrió. & RNF-011--013 y FR-006/039 de
F-03; respaldo técnico F-09. \\
\end{longtable}
}

Los actores son roles documentados con correlato de dominio parcial, no
participantes entrevistados ni nombramientos acreditados. Un rol técnico
no concede automáticamente función académica.

\section{Historias de usuario
depuradas}\label{dossier-historias-de-usuario-depuradas}

{\def\LTcaptype{none} % do not increment counter
\begin{longtable}[]{@{}
  >{\raggedright\arraybackslash}p{(\linewidth - 4\tabcolsep) * \real{0.3333}}
  >{\raggedright\arraybackslash}p{(\linewidth - 4\tabcolsep) * \real{0.3333}}
  >{\raggedright\arraybackslash}p{(\linewidth - 4\tabcolsep) * \real{0.3333}}@{}}
\toprule\noalign{}
\begin{minipage}[b]{\linewidth}\raggedright
ID / prioridad
\end{minipage} & \begin{minipage}[b]{\linewidth}\raggedright
Historia
\end{minipage} & \begin{minipage}[b]{\linewidth}\raggedright
Origen y mejora
\end{minipage} \\
\midrule\noalign{}
\endhead
\bottomrule\noalign{}
\endlastfoot
HU-AA-01 / P1 & Como usuario habilitado, quiero iniciar y terminar mi
acceso y consultar mis funciones vigentes, para operar mientras mi
sesión esté autorizada y finalizarla al concluir. & APE2/US1; conserva
beneficio y distingue sesión de permisos. \\
HU-AA-02 / P1 & Como administrador funcional autorizado, quiero
registrar unidades y carreras con datos válidos, para disponer de un
catálogo identificable antes de preparar evaluaciones. & Alta de
APE2/US2; el registro de unidad puede ser prerrequisito con fixture. \\
HU-AA-03 / P1 & Como lector autorizado, quiero consultar catálogo,
detalle e historial visibles, para identificar información de mi ámbito
sin acceder a datos ajenos. & Consulta separada de APE2/US2; no concede
modificación. \\
HU-AA-04 / P2 & Como administrador funcional autorizado, quiero corregir
datos de una carrera conservando revisiones, para actualizar el catálogo
sin perder el estado anterior ni sobrescribir cambios ajenos. & Edición
separada de APE2/US2; criterio propio de concurrencia. \\
HU-AA-05 / P2 & Como administrador funcional autorizado, quiero
desactivar una carrera con motivo, para impedir su uso en nuevos
procesos conservando identidad y referencias históricas. & Cambio de
estado separado de APE2/US2; no se presupone cierre de procesos
anteriores. \\
HU-AA-06 / P2 & Como coordinador autorizado para gestionar procesos,
quiero preparar y consultar una autoevaluación con carrera, período y
modelo publicado, para identificarla antes de tramitar su autorización
institucional. & APE2/US3; no incluye activación ni presupone autoridad
de inicio. \\
\end{longtable}
}

La división de US2 no aumenta el alcance; hace revisables sus objetivos
de alta, lectura, corrección y desactivación. Cada historia puede
comprobarse con prerrequisitos sintéticos existentes.

\section{Requisitos funcionales}\label{dossier-requisitos-funcionales}

Los IDs AA-RF son locales al dossier. La columna de origen conserva el
enlace a APE2 y a los FR del producto; no los renumera oficialmente.

{\def\LTcaptype{none} % do not increment counter
\begin{longtable}[]{@{}
  >{\raggedright\arraybackslash}p{(\linewidth - 4\tabcolsep) * \real{0.3333}}
  >{\raggedright\arraybackslash}p{(\linewidth - 4\tabcolsep) * \real{0.3333}}
  >{\raggedright\arraybackslash}p{(\linewidth - 4\tabcolsep) * \real{0.3333}}@{}}
\toprule\noalign{}
\begin{minipage}[b]{\linewidth}\raggedright
ID
\end{minipage} & \begin{minipage}[b]{\linewidth}\raggedright
Requisito observable depurado
\end{minipage} & \begin{minipage}[b]{\linewidth}\raggedright
Origen
\end{minipage} \\
\midrule\noalign{}
\endhead
\bottomrule\noalign{}
\endlastfoot
AA-RF-01 & Permitir iniciar sesión a una identidad habilitada con
credenciales válidas y rechazar credenciales inválidas sin confirmar si
la identidad existe. & APE-FR-01; FR-001/005. \\
AA-RF-02 & Permitir cerrar la sesión actual y rechazar su reutilización
posterior para operaciones protegidas. & APE-FR-01; FR-001. \\
AA-RF-03 & Permitir a la identidad autenticada consultar sus funciones y
ámbitos vigentes, sin devolver los de otra identidad. & APE-FR-01;
FR-001/002. \\
AA-RF-04 & Comprobar capacidad y ámbito vigentes antes de cada operación
protegida; leer no concederá modificar y un recurso ajeno no expondrá
información protegida. & APE-FR-02; FR-002/003. \\
AA-RF-05 & Permitir al administrador autorizado registrar una unidad con
código normalizado único y nombre, conservando identidad y estado del
alta. & Alta de unidad de APE-FR-03; FR-043. \\
AA-RF-06 & Permitir registrar una carrera con código normalizado único,
nombre, unidad activa y vigencia coherente; rechazar entrada inválida
sin crear una carrera parcial. & Alta de carrera de APE-FR-03;
FR-043. \\
AA-RF-07 & Permitir consultar catálogo, detalle e historial de
carreras/unidades según la operación, filtrando resultados y totales por
la visibilidad efectiva del actor. & APE-FR-04/05; FR-043/002. \\
AA-RF-08 & Permitir modificar datos válidos de carrera conservando
identidad, nueva revisión e historial anterior; no completar el cambio
sobre una revisión obsoleta. & APE-FR-05; FR-043; RNF-011. \\
AA-RF-09 & Permitir desactivar una carrera con motivo y revisión vigente
sin eliminación física; no activar nuevos procesos sobre una carrera no
utilizable. & APE-FR-05/07; FR-043/011. \\
AA-RF-10 & Permitir consultar e identificar la versión publicada del
modelo disponible para preparar el proceso, sin sustituir
automáticamente la elegida por otra versión. & APE-FR-06; FR-010/012. \\
AA-RF-11 & Permitir al actor con gestión vigente crear un proceso para
carrera/unidad activas, período permitido y versión publicada; guardarlo
en preparación sin activación implícita. & APE-FR-07; FR-011. \\
AA-RF-12 & Permitir al actor autorizado consultar el proceso preparado y
recuperar sus referencias y estado; rechazar una consulta ajena sin
exponer el recurso. & APE-FR-08; FR-002/011/012. \\
AA-RF-13 & Registrar cambios críticos de catálogo y creación de procesos
con actor, fecha, acción, objeto y resultado, conforme al catálogo de
auditoría aplicable; no declarar completado un cambio cuya auditoría
obligatoria no se conserve. & APE-FR-09; FR-006 y refinamiento
FR-006a. \\
AA-RF-14 & Rechazar solicitudes inválidas, conflictos o falta de acceso
con estructura y código estables que permitan identificar el problema
sin detalles técnicos sensibles. & APE-FR-10; FR-039. \\
\end{longtable}
}

No se establece una revisión inicial uniforme entre agregados ni un
código HTTP único para todas las validaciones. Los contratos concretos
de APE2 conservan esas diferencias.

\section{Requisitos no funcionales
retenidos}\label{dossier-requisitos-no-funcionales-retenidos}

{\def\LTcaptype{none} % do not increment counter
\begin{longtable}[]{@{}
  >{\raggedright\arraybackslash}p{(\linewidth - 4\tabcolsep) * \real{0.3333}}
  >{\raggedright\arraybackslash}p{(\linewidth - 4\tabcolsep) * \real{0.3333}}
  >{\raggedright\arraybackslash}p{(\linewidth - 4\tabcolsep) * \real{0.3333}}@{}}
\toprule\noalign{}
\begin{minipage}[b]{\linewidth}\raggedright
ID
\end{minipage} & \begin{minipage}[b]{\linewidth}\raggedright
Atributo y condición verificable
\end{minipage} & \begin{minipage}[b]{\linewidth}\raggedright
Fuente/origen y estado
\end{minipage} \\
\midrule\noalign{}
\endhead
\bottomrule\noalign{}
\endlastfoot
AA-RNF-01 & Contraseñas mediante hash adaptativo con sal, sin
almacenamiento reversible ni exposición en respuestas. & RNF-001, parte
de APE-NFR-01; revisión de controles, no certificación. \\
AA-RNF-02 & La sesión expirará tras 30 minutos sin actividad y la
desactivación de identidad invalidará su acceso posterior. & RNF-002,
parte de APE-NFR-01; umbral del producto F-03. \\
AA-RNF-03 & Respuestas y registros de los recorridos seleccionados no
expondrán contraseñas, tokens completos ni datos ajenos al ámbito
permitido. & RNF-005; privacidad separada de login y de estructura de
error. \\
AA-RNF-04 & El p95 de resultado visible de las operaciones interactivas
seleccionadas será \ensuremath{\leq} 2 s bajo el entorno/carga piloto que se acuerde;
archivos e informes están excluidos. & RNF-006, APE-NFR-02; en revisión
hasta fijar protocolo. \\
AA-RNF-05 & Una actualización concurrente basada en revisión obsoleta
será detectada antes de sobrescribir y conservará el cambio exitoso
anterior. & RNF-011, APE-NFR-04; propiedad transversal complementaria de
AA-RF-08. \\
AA-RNF-06 & La auditoría será inmutable para usuarios ordinarios y
consultable únicamente con permisos aplicables. & RNF-012, parte de
APE-NFR-05; retención institucional pendiente. \\
AA-RNF-07 & Cada solicitud tendrá correlación que permita relacionar
respuesta y diagnóstico técnico sin revelar datos sensibles. & RNF-013,
parte de APE-NFR-05; comprobar errores y éxitos. \\
AA-RNF-08 & Los recorridos seleccionados se completarán con teclado,
foco visible y controles identificables; el objetivo de accesibilidad es
WCAG 2.2 AA con revisión de todos los criterios aplicables. & RNF-014,
APE-NFR-06; F-13; no certificado. \\
AA-RNF-09 & Las reglas de estado/autorización estarán separadas de
controladores y persistencia y tendrán comprobaciones reproducibles con
datos sintéticos. & RNF-017/018, APE-NFR-07; revisión de código y
dependencias. \\
AA-RNF-10 & Los intercambios HTTP seleccionados tendrán método, entrada,
representación, cabeceras, resultado y errores documentados y coherentes
con el contrato aceptado. & APE-NFR-08; F-09/F-14 y FR-039; revisión
técnica abierta HTTP-AA-01. \\
\end{longtable}
}

\textbf{Decisiones de depuración:} APE-NFR-10 (datos sintéticos) se
conserva como restricción del trabajo y de todas las pruebas, evitando
presentarlo además como atributo duplicado del comportamiento
productivo. APE-NFR-03 (canales complementarios) se mantiene en backlog
global RNF-009: el recorte no contiene un canal operativo que permita
justificar su ensayo aquí. APE-NFR-09 (19 de 20 participantes) sigue
como objetivo global SC-007 y pregunta metodológica; no se promete un
estudio humano dentro del AA individual. No se elimina ninguno del
producto ni se afirma aprobación docente del aplazamiento.

\section{Criterios de aceptación
funcionales}\label{dossier-criterios-de-aceptaciuxf3n-funcionales}

Los AC-AA reemplazan la agrupación local de criterios APE2 solo en este
dossier. Los casos de entrada inválida se separan para que su ejecución
posterior sea interpretable.

{\def\LTcaptype{none} % do not increment counter
\begin{longtable}[]{@{}
  >{\raggedright\arraybackslash}p{(\linewidth - 2\tabcolsep) * \real{0.2500}}
  >{\raggedright\arraybackslash}p{(\linewidth - 2\tabcolsep) * \real{0.7500}}@{}}
\toprule\noalign{}
\begin{minipage}[b]{\linewidth}\raggedright
ID
\end{minipage} & \begin{minipage}[b]{\linewidth}\raggedright
Criterio Dado--Cuando--Entonces
\end{minipage} \\
\midrule\noalign{}
\endhead
\bottomrule\noalign{}
\endlastfoot
AC-AA-01 & \textbf{Dado que} la identidad está habilitada y las
credenciales son válidas, \textbf{cuando} inicia sesión,
\textbf{entonces} obtiene una sesión para esa identidad. \\
AC-AA-02 & \textbf{Dado que} la identidad no existe o la contraseña es
incorrecta, \textbf{cuando} intenta iniciar sesión, \textbf{entonces} no
se concede sesión ni se revela cuál de las dos causas ocurrió. \\
AC-AA-03 & \textbf{Dado que} existe una sesión activa, \textbf{cuando}
se cierra y se reutiliza, \textbf{entonces} la solicitud protegida es
rechazada. \\
AC-AA-04 & \textbf{Dado que} hay sesión vigente, \textbf{cuando}
consulta su identidad, \textbf{entonces} obtiene sus funciones/ámbitos
vigentes y no los de otra persona. \\
AC-AA-05 & \textbf{Dado que} no hay sesión válida, \textbf{cuando} se
pide una operación protegida, \textbf{entonces} no se entrega su
representación. \\
AC-AA-06 & \textbf{Dado que} una identidad carece de capacidad de
administrar, \textbf{cuando} intenta alta o cambio de catálogo,
\textbf{entonces} se rechaza y el catálogo permanece igual. \\
AC-AA-07 & \textbf{Dado que} existe una unidad con código normalizado,
\textbf{cuando} se intenta registrar otro equivalente, \textbf{entonces}
se rechaza sin duplicar la unidad. \\
AC-AA-08 & \textbf{Dado que} la unidad está activa y el código de
carrera es nuevo, \textbf{cuando} el administrador registra datos
válidos, \textbf{entonces} se conserva la carrera con unidad y
vigencia. \\
AC-AA-09 & \textbf{Dado que} ya existe el código normalizado de carrera,
\textbf{cuando} se envía un alta equivalente, \textbf{entonces} se
rechaza sin crear otra carrera. \\
AC-AA-10 & \textbf{Dado que} falta un campo obligatorio o un dato tiene
formato inválido, \textbf{cuando} se solicita el alta, \textbf{entonces}
se informa el error y no se crea el recurso. \\
AC-AA-11 & \textbf{Dado que} la unidad no existe o no está activa,
\textbf{cuando} se registra carrera, \textbf{entonces} no se crea la
carrera. \\
AC-AA-12 & \textbf{Dado que} la vigencia termina antes de su inicio,
\textbf{cuando} se registra carrera, \textbf{entonces} se rechaza y no
queda una carrera parcial. \\
AC-AA-13 & \textbf{Dado que} el lector solo tiene carrera A autorizada,
\textbf{cuando} consulta lista, total o detalle de B, \textbf{entonces}
no recibe datos ni totales ajenos a su visibilidad. \\
AC-AA-14 & \textbf{Dado que} la revisión está vigente y los nuevos datos
son válidos, \textbf{cuando} modifica la carrera, \textbf{entonces}
cambia su revisión conservando identidad y estado anterior
consultable. \\
AC-AA-15 & \textbf{Dado que} dos clientes parten de la misma revisión y
uno guarda, \textbf{cuando} el otro envía la revisión anterior,
\textbf{entonces} su cambio se rechaza y el primero permanece
intacto. \\
AC-AA-16 & \textbf{Dado que} la carrera tiene historial, \textbf{cuando}
se desactiva con motivo válido y revisión vigente, \textbf{entonces}
conserva identidad e historial y no es utilizable para un proceso
nuevo. \\
AC-AA-17 & \textbf{Dado que} no existe un cambio de estado,
\textbf{cuando} se solicita repetirlo, \textbf{entonces} no se crea una
transición ficticia. \\
AC-AA-18 & \textbf{Dado que} el actor tiene acceso al modelo,
\textbf{cuando} consulta una versión publicada identificada,
\textbf{entonces} recupera su identidad/número/estado sin reemplazo
automático por otra. \\
AC-AA-19 & \textbf{Dado que} hay gestión autorizada, carrera/unidad
activas, período válido y versión publicada, \textbf{cuando} se prepara
el proceso, \textbf{entonces} queda PREPARATION con las referencias
elegidas. \\
AC-AA-20 & \textbf{Dado que} la carrera está desactivada o no es
utilizable, \textbf{cuando} se intenta preparar un proceso nuevo,
\textbf{entonces} no se crea el proceso. \\
AC-AA-21 & \textbf{Dado que} el período es inválido o excede la vigencia
permitida, \textbf{cuando} se prepara el proceso, \textbf{entonces} se
rechaza sin persistencia parcial. \\
AC-AA-22 & \textbf{Dado que} la versión no está publicada/disponible,
\textbf{cuando} se intenta preparar proceso, \textbf{entonces} no se
crea ni se elige silenciosamente otra versión. \\
AC-AA-23 & \textbf{Dado que} el actor tiene acceso al proceso preparado,
\textbf{cuando} consulta, \textbf{entonces} recupera carrera, período,
modelo y estado de preparación sin acreditación inferida. \\
AC-AA-24 & \textbf{Dado que} el recurso es ajeno al ámbito del lector,
\textbf{cuando} intenta acceder al proceso, \textbf{entonces} no recibe
su representación protegida. \\
AC-AA-25 & \textbf{Dado que} una escritura exige auditoría y su
persistencia falla, \textbf{cuando} se ejecuta el comando,
\textbf{entonces} no se presenta como cambio completado sin el registro
obligatorio. \\
AC-AA-26 & \textbf{Dado que} una solicitud debe rechazarse,
\textbf{cuando} se responde, \textbf{entonces} la estructura documentada
identifica código, descripción y correlación sin secretos. \\
AC-AA-27 & \textbf{Dado que} el código de unidad es nuevo y el nombre
válido, \textbf{cuando} el administrador registra la unidad,
\textbf{entonces} obtiene identidad, estado y revisión y puede
consultarla. \\
\end{longtable}
}

AC-AA-27 incorpora el éxito de alta de unidad para no depender
únicamente de su rechazo por duplicado. Se distinguen 27 AC funcionales
y diez VC de calidad; son verificaciones especificadas, no resultados
ejecutados.

\section{Verificaciones de atributos de
calidad}\label{dossier-verificaciones-de-atributos-de-calidad}

{\def\LTcaptype{none} % do not increment counter
\begin{longtable}[]{@{}
  >{\raggedright\arraybackslash}p{(\linewidth - 2\tabcolsep) * \real{0.2500}}
  >{\raggedright\arraybackslash}p{(\linewidth - 2\tabcolsep) * \real{0.7500}}@{}}
\toprule\noalign{}
\begin{minipage}[b]{\linewidth}\raggedright
ID
\end{minipage} & \begin{minipage}[b]{\linewidth}\raggedright
Verificación propuesta, sin ejecución declarada
\end{minipage} \\
\midrule\noalign{}
\endhead
\bottomrule\noalign{}
\endlastfoot
VC-AA-01 & Con credenciales sintéticas, inspeccionar hashes y comprobar
algoritmo adaptativo/sal y ausencia de contraseña reversible o en
respuestas. \\
VC-AA-02 & Con reloj controlado, comprobar solicitud antes del límite,
al cumplir 30 minutos sin actividad y después de desactivar la
identidad. El límite se mide desde última actividad válida. \\
VC-AA-03 & Ejecutar casos de éxito/rechazo con secretos de prueba y
comprobar cuerpos y registros: no deben contenerlos completos ni datos
ajenos. \\
VC-AA-04 & Acordar entorno y recorrido; medir desde acción en navegador
hasta resultado visible y calcular p95. Los 2 s vienen de F-03; la
campaña histórica F-12 no sustituye esta medición. \\
VC-AA-05 & Repetir AC-AA-15 con dos actores autorizados; una revisión
obsoleta debe causar cero sobrescrituras y conservar el cambio
confirmado. \\
VC-AA-06 & Con usuario ordinario, intentar modificar auditoría y
consultar fuera de ámbito; no debe cambiar eventos ni obtener registros
no autorizados. \\
VC-AA-07 & Comprobar correlación presente y consistente en respuestas de
éxito/error y diagnóstico autorizado correspondiente, sin secretos. \\
VC-AA-08 & Completar login, catálogo y preparación con teclado; revisar
foco, nombres, contraste y demás criterios AA aplicables; registrar
criterio, pantalla y evidencia. \\
VC-AA-09 & Revisar dependencias de reglas, controles y persistencia;
ejecutar posteriormente reglas de arquitectura y pruebas relevantes,
incluyendo revisión de SQL que ArchUnit no interpreta. \\
VC-AA-10 & Comparar contratos con respuestas reales: tipos JSON, método,
estado y cabeceras aplicables; comprobar además HTTP-AA-01 antes de
afirmar compatibilidad completa de autenticación. \\
\end{longtable}
}

\chapter{Trazabilidad, depuración y decisiones
pendientes}\label{dossier-trazabilidad-depuraciuxf3n-y-decisiones-pendientes}

\section{Matriz inicial y lectura de
estados}\label{dossier-matriz-inicial-y-lectura-de-estados}

La \href{MATRIZ_TRAZABILIDAD.md}{matriz de trazabilidad} contiene una
fila por cada AA-RF y AA-RNF retenido. Sus 24 filas vinculan actor,
necesidad, historia/requisito, AC o VC, fuente y estado. No se usa
``Validado'' para un elemento que solo tiene respaldo documental.

Las historias aportan beneficio; los requisitos definen obligación; los
AC/VC definen comprobación; las fuentes justifican origen o
implementación. Un vínculo a código es contraste, no causa suficiente
para conservar una funcionalidad innecesaria.

\section{Ejemplos de justificación
completa}\label{dossier-ejemplos-de-justificaciuxf3n-completa}

\textbf{Alta de carrera:} A-AA-02 necesita N-AA-02; HU-AA-02 y AA-RF-06
definen el alta válida; AC-AA-08--12 distinguen éxito, duplicado,
formato, unidad y fecha; F-03/F-06 justifican el catálogo y C-03/C-04
permiten contrastar su comportamiento. F-08 aporta contexto
institucional, pero no es una orden de implementar un endpoint de
carreras.

\textbf{Preparación de proceso:} A-AA-03 necesita N-AA-03; HU-AA-06 y
AA-RF-11 exigen referencias exactas y PREPARATION; AC-AA-19--22 cubren
sus condiciones. F-07/F-08 sostienen el contexto, F-03/F-04 la decisión
del producto y C-05 el contraste. Las designaciones y activación siguen
pendientes, fuera del recorte.

\textbf{Concurrencia:} A-AA-02 necesita N-AA-05; AA-RF-08 es la
operación y AA-RNF-05 su calidad transversal. AC-AA-15/VC-AA-05
comprueban pérdida de cambios; RNF-011 de F-03 es la obligación y F-09
el fundamento de la precondición HTTP. No se afirma una prueba superada
por existir ETag.

\section{Calidad y ausencia de
redundancia}\label{dossier-calidad-y-ausencia-de-redundancia}

{\def\LTcaptype{none} % do not increment counter
\begin{longtable}[]{@{}
  >{\raggedright\arraybackslash}p{(\linewidth - 4\tabcolsep) * \real{0.3333}}
  >{\raggedright\arraybackslash}p{(\linewidth - 4\tabcolsep) * \real{0.3333}}
  >{\raggedright\arraybackslash}p{(\linewidth - 4\tabcolsep) * \real{0.3333}}@{}}
\toprule\noalign{}
\begin{minipage}[b]{\linewidth}\raggedright
Aspecto
\end{minipage} & \begin{minipage}[b]{\linewidth}\raggedright
Acción de depuración
\end{minipage} & \begin{minipage}[b]{\linewidth}\raggedright
Límite restante
\end{minipage} \\
\midrule\noalign{}
\endhead
\bottomrule\noalign{}
\endlastfoot
Claridad & Separar login, logout e identidad; distinguir alta, lectura,
edición y desactivación. & Confirmar vocabulario de roles
institucionales. \\
Necesidad & Cada requisito retenido se vincula a N-AA-01--05. &
Necesidades del equipo documentadas, sin entrevista real aportada. \\
Consistencia & Mantener recorte U1 y no convertir preparación en
activación. & CL-APE-01 aún requiere aceptación. \\
Verificabilidad & Separar validaciones compuestas de APE2 en
AC-AA-10/11/12 y AC-AA-20/21/22. & Ejecución y protocolos de
carga/humanos pendientes. \\
No redundancia & Separar evento auditado de inmutabilidad; operación de
edición de concurrencia; condición de ensayo de atributo de producto. &
Conservar enlaces de refinamiento para evitar contar dos veces
cobertura. \\
Compatibilidad con fuentes & Contexto oficial y política del software
separados; HTTP-AA-01 explícito. & Corroborar vigencia, roles y
comportamiento de 401 en ejecución. \\
\end{longtable}
}

\section{Observaciones y decisiones
pendientes}\label{dossier-observaciones-y-decisiones-pendientes}

{\def\LTcaptype{none} % do not increment counter
\begin{longtable}[]{@{}
  >{\raggedright\arraybackslash}p{(\linewidth - 4\tabcolsep) * \real{0.3333}}
  >{\raggedright\arraybackslash}p{(\linewidth - 4\tabcolsep) * \real{0.3333}}
  >{\raggedright\arraybackslash}p{(\linewidth - 4\tabcolsep) * \real{0.3333}}@{}}
\toprule\noalign{}
\begin{minipage}[b]{\linewidth}\raggedright
ID
\end{minipage} & \begin{minipage}[b]{\linewidth}\raggedright
Decisión o vacío
\end{minipage} & \begin{minipage}[b]{\linewidth}\raggedright
Estado y tratamiento
\end{minipage} \\
\midrule\noalign{}
\endhead
\bottomrule\noalign{}
\endlastfoot
P-AA-01 & Confirmar recorte U1 y autoridad de cada rol. & Pendiente
docente/equipo y designaciones; hereda CL-APE-01/02. \\
P-AA-02 & Verificar vigencia y aplicabilidad institucional de F-07/F-08.
& Fuente primaria consultada; no se cierran CL-011--017 ni se reescriben
manifiestos. \\
P-AA-03 & Aportar entrevistas o diagnóstico del proceso actual. & No
inventar cifras de tiempo o pérdida documental; hereda CL-APE-04. \\
P-AA-04 & Resolver cinco fallos consecutivos frente a acumulados. &
Mantener RNF-003/CL-APE-03 como pendiente; no elegir política por código
existente. \\
P-AA-05 & Acordar entorno/carga de p95 y protocolo humano de usabilidad.
& AA-RNF-04 en revisión; SC-007 global no medido aquí. \\
P-AA-06 & Comprobar WWW-Authenticate en 401 y desafío aplicable. &
HTTP-AA-01: contraste técnico pendiente; no se inventa JWT. \\
P-AA-07 & Confirmar utilidad de ensayar canales complementarios
ausentes. & APE-NFR-03 permanece global/diferido, fuera de la selección
activa AA2. \\
P-AA-08 & Identificar estudiante, revisión personal y enlace de
contribución. & La asistencia de IA no acredita una participación
individual. \\
\end{longtable}
}

\section{Producto, revisión individual y tiempo
previsto}\label{dossier-producto-revisiuxf3n-individual-y-tiempo-previsto}

El producto local incluye dossier, matriz, fuentes, nota de cambios y
versión LaTeX. El remoto observado es
\href{https://github.com/M16artz/SIGEAC}{M16artz/SIGEAC}; no se verificó
acceso externo ni se publicó este AA2. El commit base no es un commit de
la actividad.

Distribución sugerida de 4,5 horas, no cronometraje: 45 min para RFC; 45
min para fuentes del dominio; 60 min para historias/requisitos; 60 min
para matriz/criterios; 45 min para dossier/nota; 15 min para revisión y
evidencia individual. Total: 270 minutos.

El estudiante debe revisar los cambios de la
\href{NOTA_CAMBIOS_INDIVIDUAL.md}{nota individual}, justificar con sus
palabras los elementos modificados e identificar su evidencia real. No
se atribuyen correcciones al docente ni a un integrante sin registro.

\section{Declaración de uso de IA y
límites}\label{dossier-declaraciuxf3n-de-uso-de-ia-y-luxedmites}

Se utilizó ChatGPT/Codex para revisar archivos y fuentes públicas,
depurar textos, construir la matriz y generar documentos. No se
realizaron entrevistas, validación con docente, solicitudes al backend,
pruebas de negocio, commits ni publicación. Las comprobaciones
documentales no equivalen a cumplimiento funcional o institucional. La
revisión personal, datos del estudiante y aprobación del equipo
permanecen pendientes.

\cleardoublepage
\nocite{proyecto,ape2,rfc9110,caces2024,unl2022,wcag22,rfc5789}
\bibliographystyle{IEEEtran}
\bibliography{referencias}
\appendix
\chapter{Anexos}
\section{Matriz inicial de trazabilidad --- fichas
completas}\label{matriz-matriz-inicial-de-trazabilidad-fichas-completas}

Los 24 registros conservan todos los campos de la matriz Markdown. TR
significa trazado documental; DP, confirmación de dominio pendiente; VM,
verificación o medición pendiente; RT, revisión técnica pendiente. No
son aprobaciones ni resultados de ejecución.

\subsection{AA-RF-01}\label{matriz-aa-rf-01}

{\def\LTcaptype{none} % do not increment counter
\begin{longtable}[]{@{}
  >{\raggedright\arraybackslash}p{(\linewidth - 2\tabcolsep) * \real{0.2500}}
  >{\raggedright\arraybackslash}p{(\linewidth - 2\tabcolsep) * \real{0.7500}}@{}}
\toprule\noalign{}
\begin{minipage}[b]{\linewidth}\raggedright
Campo
\end{minipage} & \begin{minipage}[b]{\linewidth}\raggedright
Relación
\end{minipage} \\
\midrule\noalign{}
\endhead
\bottomrule\noalign{}
\endlastfoot
Actor & A-AA-01, usuario habilitado \\
Necesidad & N-AA-01: acceder sin exponer identidad ajena. \\
Historia / requisito & HU-AA-01: conceder sesión solo a identidad
habilitada con credenciales válidas. \\
Criterio de aceptación o verificación & AC-AA-01/02: Dado acceso
válido/erróneo, cuando autentica, entonces concede solo el válido sin
enumerar cuentas. \\
Fuente & F-02 APE-FR-01; F-03 FR-001/005; F-11 C-01/02. \\
Estado & TR; VM de autenticación. \\
\end{longtable}
}

\subsection{AA-RF-02}\label{matriz-aa-rf-02}

{\def\LTcaptype{none} % do not increment counter
\begin{longtable}[]{@{}
  >{\raggedright\arraybackslash}p{(\linewidth - 2\tabcolsep) * \real{0.2500}}
  >{\raggedright\arraybackslash}p{(\linewidth - 2\tabcolsep) * \real{0.7500}}@{}}
\toprule\noalign{}
\begin{minipage}[b]{\linewidth}\raggedright
Campo
\end{minipage} & \begin{minipage}[b]{\linewidth}\raggedright
Relación
\end{minipage} \\
\midrule\noalign{}
\endhead
\bottomrule\noalign{}
\endlastfoot
Actor & A-AA-01, usuario habilitado \\
Necesidad & N-AA-01: terminar acceso. \\
Historia / requisito & HU-AA-01: cerrar e invalidar sesión. \\
Criterio de aceptación o verificación & AC-AA-03: Dado sesión activa,
cuando cierra y reutiliza, entonces se rechaza acceso protegido. \\
Fuente & F-02 APE-FR-01; F-03 FR-001; F-11 C-01/02. \\
Estado & TR; VM de revocación. \\
\end{longtable}
}

\subsection{AA-RF-03}\label{matriz-aa-rf-03}

{\def\LTcaptype{none} % do not increment counter
\begin{longtable}[]{@{}
  >{\raggedright\arraybackslash}p{(\linewidth - 2\tabcolsep) * \real{0.2500}}
  >{\raggedright\arraybackslash}p{(\linewidth - 2\tabcolsep) * \real{0.7500}}@{}}
\toprule\noalign{}
\begin{minipage}[b]{\linewidth}\raggedright
Campo
\end{minipage} & \begin{minipage}[b]{\linewidth}\raggedright
Relación
\end{minipage} \\
\midrule\noalign{}
\endhead
\bottomrule\noalign{}
\endlastfoot
Actor & A-AA-01, usuario habilitado \\
Necesidad & N-AA-01: conocer funciones vigentes. \\
Historia / requisito & HU-AA-01: consultar identidad y ámbitos
propios. \\
Criterio de aceptación o verificación & AC-AA-04: Dado sesión vigente,
cuando consulta identidad, entonces recibe únicamente sus
funciones/ámbitos. \\
Fuente & F-03 FR-001/002; F-05 P01; F-11 C-01. \\
Estado & TR; DP del catálogo de funciones. \\
\end{longtable}
}

\subsection{AA-RF-04}\label{matriz-aa-rf-04}

{\def\LTcaptype{none} % do not increment counter
\begin{longtable}[]{@{}
  >{\raggedright\arraybackslash}p{(\linewidth - 2\tabcolsep) * \real{0.2500}}
  >{\raggedright\arraybackslash}p{(\linewidth - 2\tabcolsep) * \real{0.7500}}@{}}
\toprule\noalign{}
\begin{minipage}[b]{\linewidth}\raggedright
Campo
\end{minipage} & \begin{minipage}[b]{\linewidth}\raggedright
Relación
\end{minipage} \\
\midrule\noalign{}
\endhead
\bottomrule\noalign{}
\endlastfoot
Actor & A-AA-01/02/03/04 \\
Necesidad & N-AA-01/04: actuar/leer solo dentro de permiso. \\
Historia / requisito & HU-AA-01--06: comprobar capacidad y ámbito por
operación. \\
Criterio de aceptación o verificación & AC-AA-05/06/13/24: Dado falta de
sesión, permiso o ámbito, cuando opera, entonces no modifica ni revela
recursos ajenos. \\
Fuente & F-03 FR-002/003; F-04 responsabilidades; F-11 C-04/05. \\
Estado & TR; DP de autoridades y VM de cobertura. \\
\end{longtable}
}

\subsection{AA-RF-05}\label{matriz-aa-rf-05}

{\def\LTcaptype{none} % do not increment counter
\begin{longtable}[]{@{}
  >{\raggedright\arraybackslash}p{(\linewidth - 2\tabcolsep) * \real{0.2500}}
  >{\raggedright\arraybackslash}p{(\linewidth - 2\tabcolsep) * \real{0.7500}}@{}}
\toprule\noalign{}
\begin{minipage}[b]{\linewidth}\raggedright
Campo
\end{minipage} & \begin{minipage}[b]{\linewidth}\raggedright
Relación
\end{minipage} \\
\midrule\noalign{}
\endhead
\bottomrule\noalign{}
\endlastfoot
Actor & A-AA-02, administrador funcional \\
Necesidad & N-AA-02: identificar unidad de la carrera. \\
Historia / requisito & HU-AA-02: registrar unidad válida y única. \\
Criterio de aceptación o verificación & AC-AA-07/27: Dado código
nuevo/duplicado, cuando registra, entonces crea solo el nuevo sin
duplicados. \\
Fuente & F-02 APE-FR-03; F-06 FR-043; F-11 C-03/04. \\
Estado & TR; DP de administración. \\
\end{longtable}
}

\subsection{AA-RF-06}\label{matriz-aa-rf-06}

{\def\LTcaptype{none} % do not increment counter
\begin{longtable}[]{@{}
  >{\raggedright\arraybackslash}p{(\linewidth - 2\tabcolsep) * \real{0.2500}}
  >{\raggedright\arraybackslash}p{(\linewidth - 2\tabcolsep) * \real{0.7500}}@{}}
\toprule\noalign{}
\begin{minipage}[b]{\linewidth}\raggedright
Campo
\end{minipage} & \begin{minipage}[b]{\linewidth}\raggedright
Relación
\end{minipage} \\
\midrule\noalign{}
\endhead
\bottomrule\noalign{}
\endlastfoot
Actor & A-AA-02, administrador funcional \\
Necesidad & N-AA-02: carrera inequívoca y válida. \\
Historia / requisito & HU-AA-02: alta con código, unidad y vigencia. \\
Criterio de aceptación o verificación & AC-AA-08--12: Dado datos válidos
o inválidos, cuando registra, entonces conserva solo alta válida y
rechaza cada causa sin parcialidad. \\
Fuente & F-03/F-06 FR-043; F-05 P02; F-11 C-03/04. \\
Estado & TR; VM de variantes. \\
\end{longtable}
}

\subsection{AA-RF-07}\label{matriz-aa-rf-07}

{\def\LTcaptype{none} % do not increment counter
\begin{longtable}[]{@{}
  >{\raggedright\arraybackslash}p{(\linewidth - 2\tabcolsep) * \real{0.2500}}
  >{\raggedright\arraybackslash}p{(\linewidth - 2\tabcolsep) * \real{0.7500}}@{}}
\toprule\noalign{}
\begin{minipage}[b]{\linewidth}\raggedright
Campo
\end{minipage} & \begin{minipage}[b]{\linewidth}\raggedright
Relación
\end{minipage} \\
\midrule\noalign{}
\endhead
\bottomrule\noalign{}
\endlastfoot
Actor & A-AA-04 y A-AA-02 \\
Necesidad & N-AA-04: leer datos/historial visibles. \\
Historia / requisito & HU-AA-03: consultar catálogo, detalle e historial
autorizado. \\
Criterio de aceptación o verificación & AC-AA-13/14: Dado visibilidad
acotada e historial, cuando consulta, entonces obtiene únicamente
recursos y revisiones permitidos. \\
Fuente & F-02 APE-FR-04/05; F-03 FR-002; F-06 FR-043; C-03/04. \\
Estado & TR; VM de listados/totales. \\
\end{longtable}
}

\subsection{AA-RF-08}\label{matriz-aa-rf-08}

{\def\LTcaptype{none} % do not increment counter
\begin{longtable}[]{@{}
  >{\raggedright\arraybackslash}p{(\linewidth - 2\tabcolsep) * \real{0.2500}}
  >{\raggedright\arraybackslash}p{(\linewidth - 2\tabcolsep) * \real{0.7500}}@{}}
\toprule\noalign{}
\begin{minipage}[b]{\linewidth}\raggedright
Campo
\end{minipage} & \begin{minipage}[b]{\linewidth}\raggedright
Relación
\end{minipage} \\
\midrule\noalign{}
\endhead
\bottomrule\noalign{}
\endlastfoot
Actor & A-AA-02, administrador funcional \\
Necesidad & N-AA-02/05: corregir sin perder historia. \\
Historia / requisito & HU-AA-04: actualizar con revisión vigente y
conservar predecesora. \\
Criterio de aceptación o verificación & AC-AA-14/15: Dado revisión
actual/obsoleta, cuando cambia, entonces guarda la actual y rechaza la
obsoleta sin sobrescribir. \\
Fuente & F-02 APE-FR-05; F-06 FR-043; F-03 RNF-011; C-04/06. \\
Estado & TR; VM concurrente. \\
\end{longtable}
}

\subsection{AA-RF-09}\label{matriz-aa-rf-09}

{\def\LTcaptype{none} % do not increment counter
\begin{longtable}[]{@{}
  >{\raggedright\arraybackslash}p{(\linewidth - 2\tabcolsep) * \real{0.2500}}
  >{\raggedright\arraybackslash}p{(\linewidth - 2\tabcolsep) * \real{0.7500}}@{}}
\toprule\noalign{}
\begin{minipage}[b]{\linewidth}\raggedright
Campo
\end{minipage} & \begin{minipage}[b]{\linewidth}\raggedright
Relación
\end{minipage} \\
\midrule\noalign{}
\endhead
\bottomrule\noalign{}
\endlastfoot
Actor & A-AA-02 y A-AA-03 \\
Necesidad & N-AA-02/03: no usar carreras desactivadas. \\
Historia / requisito & HU-AA-05/06: desactivar con motivo e impedir
nuevos procesos. \\
Criterio de aceptación o verificación & AC-AA-16/17/20: Dado cambio
válido/repetido o carrera inutilizable, cuando solicita, entonces
conserva historia y bloquea alta indebida. \\
Fuente & F-06 FR-043; F-03 FR-011; F-04 estados; C-04/05. \\
Estado & TR; VM de cambio/uso posterior. \\
\end{longtable}
}

\subsection{AA-RF-10}\label{matriz-aa-rf-10}

{\def\LTcaptype{none} % do not increment counter
\begin{longtable}[]{@{}
  >{\raggedright\arraybackslash}p{(\linewidth - 2\tabcolsep) * \real{0.2500}}
  >{\raggedright\arraybackslash}p{(\linewidth - 2\tabcolsep) * \real{0.7500}}@{}}
\toprule\noalign{}
\begin{minipage}[b]{\linewidth}\raggedright
Campo
\end{minipage} & \begin{minipage}[b]{\linewidth}\raggedright
Relación
\end{minipage} \\
\midrule\noalign{}
\endhead
\bottomrule\noalign{}
\endlastfoot
Actor & A-AA-03 con lectura de modelos \\
Necesidad & N-AA-03: identificar modelo exacto. \\
Historia / requisito & HU-AA-06: consultar versión publicada sin
selección implícita. \\
Criterio de aceptación o verificación & AC-AA-18/22: Dado versión
identificada, cuando consulta/prepara, entonces conserva la disponible
publicada o rechaza sin sustitución. \\
Fuente & F-03 FR-010/012; F-07 contexto; F-10 límite sintético; F-11
C-05. \\
Estado & TR; DP de aplicabilidad normativa. \\
\end{longtable}
}

\subsection{AA-RF-11}\label{matriz-aa-rf-11}

{\def\LTcaptype{none} % do not increment counter
\begin{longtable}[]{@{}
  >{\raggedright\arraybackslash}p{(\linewidth - 2\tabcolsep) * \real{0.2500}}
  >{\raggedright\arraybackslash}p{(\linewidth - 2\tabcolsep) * \real{0.7500}}@{}}
\toprule\noalign{}
\begin{minipage}[b]{\linewidth}\raggedright
Campo
\end{minipage} & \begin{minipage}[b]{\linewidth}\raggedright
Relación
\end{minipage} \\
\midrule\noalign{}
\endhead
\bottomrule\noalign{}
\endlastfoot
Actor & A-AA-03, gestor autorizado \\
Necesidad & N-AA-03: preparar antes de autorizar. \\
Historia / requisito & HU-AA-06: crear PREPARATION con referencias
exactas. \\
Criterio de aceptación o verificación & AC-AA-19--22: Dado
prerrequisitos válidos/ausentes, cuando prepara, entonces guarda solo el
proceso válido sin activación. \\
Fuente & F-03 FR-011; F-04/F-05 P05; F-07/F-08 contexto; C-05. \\
Estado & TR; DP del recorte/roles; VM. \\
\end{longtable}
}

\subsection{AA-RF-12}\label{matriz-aa-rf-12}

{\def\LTcaptype{none} % do not increment counter
\begin{longtable}[]{@{}
  >{\raggedright\arraybackslash}p{(\linewidth - 2\tabcolsep) * \real{0.2500}}
  >{\raggedright\arraybackslash}p{(\linewidth - 2\tabcolsep) * \real{0.7500}}@{}}
\toprule\noalign{}
\begin{minipage}[b]{\linewidth}\raggedright
Campo
\end{minipage} & \begin{minipage}[b]{\linewidth}\raggedright
Relación
\end{minipage} \\
\midrule\noalign{}
\endhead
\bottomrule\noalign{}
\endlastfoot
Actor & A-AA-03; lector con acceso aplicable \\
Necesidad & N-AA-04: conocer proceso permitido. \\
Historia / requisito & HU-AA-06: consultar estado y referencias, sin
datos ajenos. \\
Criterio de aceptación o verificación & AC-AA-23/24: Dado proceso
visible/ajeno, cuando consulta, entonces recupera solo el visible y no
infiere acreditación. \\
Fuente & F-03 FR-002/011/012; F-04 alcance; F-11 C-05. \\
Estado & TR; VM de permiso del proceso. \\
\end{longtable}
}

\subsection{AA-RF-13}\label{matriz-aa-rf-13}

{\def\LTcaptype{none} % do not increment counter
\begin{longtable}[]{@{}
  >{\raggedright\arraybackslash}p{(\linewidth - 2\tabcolsep) * \real{0.2500}}
  >{\raggedright\arraybackslash}p{(\linewidth - 2\tabcolsep) * \real{0.7500}}@{}}
\toprule\noalign{}
\begin{minipage}[b]{\linewidth}\raggedright
Campo
\end{minipage} & \begin{minipage}[b]{\linewidth}\raggedright
Relación
\end{minipage} \\
\midrule\noalign{}
\endhead
\bottomrule\noalign{}
\endlastfoot
Actor & A-AA-02/03 y auditor autorizado \\
Necesidad & N-AA-05: explicar cambios críticos. \\
Historia / requisito & HU-AA-02/04/05/06: registrar actor, fecha,
acción, objeto y resultado. \\
Criterio de aceptación o verificación & AC-AA-25: Dado auditoría
obligatoria fallida, cuando ejecuta, entonces no declara cambio
completado sin registro. \\
Fuente & F-03 FR-006; specs/001 FR-006a; F-11 C-04/05; F-12. \\
Estado & TR; VM de transacción y fallos. \\
\end{longtable}
}

\subsection{AA-RF-14}\label{matriz-aa-rf-14}

{\def\LTcaptype{none} % do not increment counter
\begin{longtable}[]{@{}
  >{\raggedright\arraybackslash}p{(\linewidth - 2\tabcolsep) * \real{0.2500}}
  >{\raggedright\arraybackslash}p{(\linewidth - 2\tabcolsep) * \real{0.7500}}@{}}
\toprule\noalign{}
\begin{minipage}[b]{\linewidth}\raggedright
Campo
\end{minipage} & \begin{minipage}[b]{\linewidth}\raggedright
Relación
\end{minipage} \\
\midrule\noalign{}
\endhead
\bottomrule\noalign{}
\endlastfoot
Actor & A-AA-01/02/03/04 \\
Necesidad & N-AA-05: comprender el rechazo. \\
Historia / requisito & HU-AA-01--06: error estable sin secretos. \\
Criterio de aceptación o verificación & AC-AA-26: Dado rechazo, cuando
responde, entonces informa código, descripción y correlación segura. \\
Fuente & F-03 FR-039; F-11 C-07/08/09; contratos APE2. \\
Estado & TR; RT HTTP-AA-01 y VM. \\
\end{longtable}
}

\subsection{AA-RNF-01}\label{matriz-aa-rnf-01}

{\def\LTcaptype{none} % do not increment counter
\begin{longtable}[]{@{}
  >{\raggedright\arraybackslash}p{(\linewidth - 2\tabcolsep) * \real{0.2500}}
  >{\raggedright\arraybackslash}p{(\linewidth - 2\tabcolsep) * \real{0.7500}}@{}}
\toprule\noalign{}
\begin{minipage}[b]{\linewidth}\raggedright
Campo
\end{minipage} & \begin{minipage}[b]{\linewidth}\raggedright
Relación
\end{minipage} \\
\midrule\noalign{}
\endhead
\bottomrule\noalign{}
\endlastfoot
Actor & A-AA-01; administrador técnico autorizado \\
Necesidad & N-AA-01: proteger credenciales. \\
Historia / requisito & HU-AA-01: hash adaptativo con sal; sin
reversibilidad. \\
Criterio de aceptación o verificación & VC-AA-01: Dado credenciales
sintéticas, cuando se inspecciona almacenamiento/respuesta, entonces no
hay secreto reversible ni expuesto. \\
Fuente & F-03 RNF-001; F-02 APE-NFR-01; F-11 C-02. \\
Estado & TR; VM de almacenamiento. \\
\end{longtable}
}

\subsection{AA-RNF-02}\label{matriz-aa-rnf-02}

{\def\LTcaptype{none} % do not increment counter
\begin{longtable}[]{@{}
  >{\raggedright\arraybackslash}p{(\linewidth - 2\tabcolsep) * \real{0.2500}}
  >{\raggedright\arraybackslash}p{(\linewidth - 2\tabcolsep) * \real{0.7500}}@{}}
\toprule\noalign{}
\begin{minipage}[b]{\linewidth}\raggedright
Campo
\end{minipage} & \begin{minipage}[b]{\linewidth}\raggedright
Relación
\end{minipage} \\
\midrule\noalign{}
\endhead
\bottomrule\noalign{}
\endlastfoot
Actor & A-AA-01, usuario habilitado \\
Necesidad & N-AA-01: limitar acceso inactivo. \\
Historia / requisito & HU-AA-01: expiración de 30 min e invalidación. \\
Criterio de aceptación o verificación & VC-AA-02: Dado reloj controlado
y límite alcanzado, cuando solicita, entonces se rechaza sesión
expirada/desactivada. \\
Fuente & F-03 RNF-002; F-02 APE-NFR-01; F-11 C-02. \\
Estado & TR; VM de frontera temporal. \\
\end{longtable}
}

\subsection{AA-RNF-03}\label{matriz-aa-rnf-03}

{\def\LTcaptype{none} % do not increment counter
\begin{longtable}[]{@{}
  >{\raggedright\arraybackslash}p{(\linewidth - 2\tabcolsep) * \real{0.2500}}
  >{\raggedright\arraybackslash}p{(\linewidth - 2\tabcolsep) * \real{0.7500}}@{}}
\toprule\noalign{}
\begin{minipage}[b]{\linewidth}\raggedright
Campo
\end{minipage} & \begin{minipage}[b]{\linewidth}\raggedright
Relación
\end{minipage} \\
\midrule\noalign{}
\endhead
\bottomrule\noalign{}
\endlastfoot
Actor & A-AA-01/02/03/04 \\
Necesidad & N-AA-01/04: privacidad. \\
Historia / requisito & HU-AA-01--06: no exponer secretos ni datos
ajenos. \\
Criterio de aceptación o verificación & VC-AA-03: Dado marcadores
secretos sintéticos, cuando consulta/inspecciona registros, entonces no
aparecen completos ni datos ajenos. \\
Fuente & F-03 RNF-005; F-09 §17 como contexto; F-11. \\
Estado & TR; VM de cobertura. \\
\end{longtable}
}

\subsection{AA-RNF-04}\label{matriz-aa-rnf-04}

{\def\LTcaptype{none} % do not increment counter
\begin{longtable}[]{@{}
  >{\raggedright\arraybackslash}p{(\linewidth - 2\tabcolsep) * \real{0.2500}}
  >{\raggedright\arraybackslash}p{(\linewidth - 2\tabcolsep) * \real{0.7500}}@{}}
\toprule\noalign{}
\begin{minipage}[b]{\linewidth}\raggedright
Campo
\end{minipage} & \begin{minipage}[b]{\linewidth}\raggedright
Relación
\end{minipage} \\
\midrule\noalign{}
\endhead
\bottomrule\noalign{}
\endlastfoot
Actor & A-AA-01/02/03/04 \\
Necesidad & N-AA-04: resultado oportuno. \\
Historia / requisito & HU-AA-01--06: p95 visible \ensuremath{\leq} 2 s bajo protocolo
acordado. \\
Criterio de aceptación o verificación & VC-AA-04: Dado entorno/carga
definidos, cuando mide resultado visible, entonces p95 satisface el
umbral del producto. \\
Fuente & F-03 RNF-006; F-02 APE-NFR-02; F-12 alcance limitado. \\
Estado & VM; entorno/carga pendientes. \\
\end{longtable}
}

\subsection{AA-RNF-05}\label{matriz-aa-rnf-05}

{\def\LTcaptype{none} % do not increment counter
\begin{longtable}[]{@{}
  >{\raggedright\arraybackslash}p{(\linewidth - 2\tabcolsep) * \real{0.2500}}
  >{\raggedright\arraybackslash}p{(\linewidth - 2\tabcolsep) * \real{0.7500}}@{}}
\toprule\noalign{}
\begin{minipage}[b]{\linewidth}\raggedright
Campo
\end{minipage} & \begin{minipage}[b]{\linewidth}\raggedright
Relación
\end{minipage} \\
\midrule\noalign{}
\endhead
\bottomrule\noalign{}
\endlastfoot
Actor & A-AA-02/03 que modifican \\
Necesidad & N-AA-05: evitar pérdida silenciosa. \\
Historia / requisito & HU-AA-04/05: precondición de revisión. \\
Criterio de aceptación o verificación & VC-AA-05 / AC-AA-15: Dado dos
clientes, cuando uno utiliza revisión vieja, entonces no sobrescribe el
cambio confirmado. \\
Fuente & F-03 RNF-011; F-09 §13.1.1; F-11 C-06. \\
Estado & TR; VM de concurrencia. \\
\end{longtable}
}

\subsection{AA-RNF-06}\label{matriz-aa-rnf-06}

{\def\LTcaptype{none} % do not increment counter
\begin{longtable}[]{@{}
  >{\raggedright\arraybackslash}p{(\linewidth - 2\tabcolsep) * \real{0.2500}}
  >{\raggedright\arraybackslash}p{(\linewidth - 2\tabcolsep) * \real{0.7500}}@{}}
\toprule\noalign{}
\begin{minipage}[b]{\linewidth}\raggedright
Campo
\end{minipage} & \begin{minipage}[b]{\linewidth}\raggedright
Relación
\end{minipage} \\
\midrule\noalign{}
\endhead
\bottomrule\noalign{}
\endlastfoot
Actor & Auditor autorizado; actores A-AA-02/03 \\
Necesidad & N-AA-05: historial confiable. \\
Historia / requisito & Historias con cambios: auditoría inmutable y
consulta autorizada. \\
Criterio de aceptación o verificación & VC-AA-06: Dado usuario
ordinario/ajeno, cuando intenta alterar/leer, entonces no modifica
eventos ni recibe los no autorizados. \\
Fuente & F-03 RNF-012; F-04 responsabilidades; F-12. \\
Estado & TR; DP de retención y VM. \\
\end{longtable}
}

\subsection{AA-RNF-07}\label{matriz-aa-rnf-07}

{\def\LTcaptype{none} % do not increment counter
\begin{longtable}[]{@{}
  >{\raggedright\arraybackslash}p{(\linewidth - 2\tabcolsep) * \real{0.2500}}
  >{\raggedright\arraybackslash}p{(\linewidth - 2\tabcolsep) * \real{0.7500}}@{}}
\toprule\noalign{}
\begin{minipage}[b]{\linewidth}\raggedright
Campo
\end{minipage} & \begin{minipage}[b]{\linewidth}\raggedright
Relación
\end{minipage} \\
\midrule\noalign{}
\endhead
\bottomrule\noalign{}
\endlastfoot
Actor & A-AA-01/02/03/04; soporte autorizado \\
Necesidad & N-AA-05: investigar fallos. \\
Historia / requisito & Historias transversales: correlación segura. \\
Criterio de aceptación o verificación & VC-AA-07: Dado éxito/error,
cuando relaciona respuesta y diagnóstico, entonces hay correlación
consistente sin secretos. \\
Fuente & F-03 RNF-013; F-11 C-07/08/09. \\
Estado & TR; VM de respuesta/registro. \\
\end{longtable}
}

\subsection{AA-RNF-08}\label{matriz-aa-rnf-08}

{\def\LTcaptype{none} % do not increment counter
\begin{longtable}[]{@{}
  >{\raggedright\arraybackslash}p{(\linewidth - 2\tabcolsep) * \real{0.2500}}
  >{\raggedright\arraybackslash}p{(\linewidth - 2\tabcolsep) * \real{0.7500}}@{}}
\toprule\noalign{}
\begin{minipage}[b]{\linewidth}\raggedright
Campo
\end{minipage} & \begin{minipage}[b]{\linewidth}\raggedright
Relación
\end{minipage} \\
\midrule\noalign{}
\endhead
\bottomrule\noalign{}
\endlastfoot
Actor & A-AA-01/02/03/04 \\
Necesidad & N-AA-01/04: completar recorridos accesibles. \\
Historia / requisito & HU-AA-01--06: teclado, foco y criterios AA
aplicables. \\
Criterio de aceptación o verificación & VC-AA-08: Dado interacción con
teclado, cuando completa recorridos y auditoría, entonces satisface
criterios aplicables con evidencia. \\
Fuente & F-03 RNF-014; F-13; F-12 con límites. \\
Estado & VM; sin certificación WCAG. \\
\end{longtable}
}

\subsection{AA-RNF-09}\label{matriz-aa-rnf-09}

{\def\LTcaptype{none} % do not increment counter
\begin{longtable}[]{@{}
  >{\raggedright\arraybackslash}p{(\linewidth - 2\tabcolsep) * \real{0.2500}}
  >{\raggedright\arraybackslash}p{(\linewidth - 2\tabcolsep) * \real{0.7500}}@{}}
\toprule\noalign{}
\begin{minipage}[b]{\linewidth}\raggedright
Campo
\end{minipage} & \begin{minipage}[b]{\linewidth}\raggedright
Relación
\end{minipage} \\
\midrule\noalign{}
\endhead
\bottomrule\noalign{}
\endlastfoot
Actor & A-AA-01--04 como beneficiarios; verificación técnica \\
Necesidad & N-AA-05: mantener reglas verificables. \\
Historia / requisito & Reglas de todas las historias separadas y
reproducibles. \\
Criterio de aceptación o verificación & VC-AA-09: Dado código y
fixtures, cuando revisa dependencias y prueba reglas, entonces detecta
acoplamientos y verifica invariantes. \\
Fuente & F-03 RNF-017/018; constitución; F-11/F-12. \\
Estado & TR; VM y revisión SQL. \\
\end{longtable}
}

\subsection{AA-RNF-10}\label{matriz-aa-rnf-10}

{\def\LTcaptype{none} % do not increment counter
\begin{longtable}[]{@{}
  >{\raggedright\arraybackslash}p{(\linewidth - 2\tabcolsep) * \real{0.2500}}
  >{\raggedright\arraybackslash}p{(\linewidth - 2\tabcolsep) * \real{0.7500}}@{}}
\toprule\noalign{}
\begin{minipage}[b]{\linewidth}\raggedright
Campo
\end{minipage} & \begin{minipage}[b]{\linewidth}\raggedright
Relación
\end{minipage} \\
\midrule\noalign{}
\endhead
\bottomrule\noalign{}
\endlastfoot
Actor & Cliente/actor A-AA-01--04 \\
Necesidad & N-AA-04/05: interpretar intercambio. \\
Historia / requisito & HU-AA-01--06: contrato HTTP coherente. \\
Criterio de aceptación o verificación & VC-AA-10: Dado contrato
aceptado, cuando compara intercambios, entonces tipos, estados y
cabeceras coinciden; comprobar HTTP-AA-01. \\
Fuente & F-02 APE-NFR-08; F-09/F-14; F-11 C-01/03/06/07/08. \\
Estado & RT; pendiente compatibilidad de 401. \\
\end{longtable}
}
\section{Registro de fuentes y
referencias}\label{fuentes-registro-de-fuentes-y-referencias}

\textbf{Fecha de consulta documental:} 2026-10-31. Las rutas locales son
relativas a la raíz del repositorio. El texto identifica qué se consultó
y qué permite justificar; no equipara consulta con vigencia
institucional ni con ejecución de pruebas.

\subsection{Fuentes del proyecto y de la
actividad}\label{fuentes-fuentes-del-proyecto-y-de-la-actividad}

{\def\LTcaptype{none} % do not increment counter
\begin{longtable}[]{@{}
  >{\raggedright\arraybackslash}p{(\linewidth - 6\tabcolsep) * \real{0.2500}}
  >{\raggedright\arraybackslash}p{(\linewidth - 6\tabcolsep) * \real{0.2500}}
  >{\raggedright\arraybackslash}p{(\linewidth - 6\tabcolsep) * \real{0.2500}}
  >{\raggedright\arraybackslash}p{(\linewidth - 6\tabcolsep) * \real{0.2500}}@{}}
\toprule\noalign{}
\begin{minipage}[b]{\linewidth}\raggedright
ID
\end{minipage} & \begin{minipage}[b]{\linewidth}\raggedright
Fuente identificable
\end{minipage} & \begin{minipage}[b]{\linewidth}\raggedright
Localizador consultado
\end{minipage} & \begin{minipage}[b]{\linewidth}\raggedright
Relación y límite
\end{minipage} \\
\midrule\noalign{}
\endhead
\bottomrule\noalign{}
\endlastfoot
F-01 & Guía AA2 proporcionada por el usuario & Actividades 1--7 y
producto esperado & Define el dossier individual, revisión de fuentes y
cadena actor--necesidad--requisito--aceptación--fuente. No es aprobación
del producto. \\
F-02 & Especificación APE2 &
\path|documentos-academicos/APE2-semana-2/specification.md|: US1--3,
APE-FR-01--10, APE-NFR-01--10 y CL-APE & Línea base académica para
depuración; sigue pendiente de validación docente. \\
F-03 & Especificación general SEC-UNL &
\path|docs/requirements/spec.md|: Contexto, Actores, FR-001--040,
RNF-001--019, SC y CL & Fuente de requisitos del equipo. Los 30 minutos
y p95 de 2 s provienen de este documento, no de RFC 9110. \\
F-04 & Reglas del dominio & \path|docs/requirements/domain-rules.md|:
responsabilidades, estados, decisiones y procedencia & Fuente de
política del software y límites institucionales; no constituye normativa
oficial por sí sola. \\
F-05 & Procesos del proyecto &
\path|docs/requirements/process-flows.md|: P01, P02 y P05 & Escenarios y
actores del incremento. \\
F-06 & Especificación reconciliada &
\path|specs/005-complete-evaluation-workflows/spec.md|: FR-043, FR-044,
FR-048, FR-067--070 & Catálogo, procedencia y separación entre
publicación, autorización y alcance. \\
F-10 & Manifiestos normativos sintéticos &
\path|tests/fixtures/models/normative-verification.json| y
\path|tests/fixtures/models/normative-manifest.json| & Declaran
\path|synthetic: true|, \path|institutionalValidity: false| y
verificaciones pendientes. No justifican vigencia ni equivalencia de los
cinco modelos. \\
F-11 & Fuentes ejecutables & Archivos C-01--C-10 de la sección 4 &
Permiten comprobar DTO y controles observados; el código no sustituye la
necesidad del actor ni prueba corrección normativa. \\
F-12 & Pruebas y resultados históricos &
\path|docs/testing/traceability.md|;
\path|tests/acceptance/results/performance-p01-p16.md|; pruebas
indicadas en el dossier & Orientan verificaciones posteriores. No se
reejecutaron y su cobertura no se extiende automáticamente al recorte
AA2. \\
\end{longtable}
}

\subsection{Fuentes primarias del dominio
consultadas}\label{fuentes-fuentes-primarias-del-dominio-consultadas}

\subsubsection{F-07 --- Modelo Genérico CACES, marzo de
2024}\label{fuentes-f-07-modelo-genuxe9rico-caces-marzo-de-2024}

\begin{itemize}
\tightlist
\item
  Autor institucional: Consejo de Aseguramiento de la Calidad de la
  Educación Superior (CACES).
\item
  Identificación: Modelo genérico para la evaluación del entorno de
  aprendizaje de carreras de grado, marzo de 2024.
\item
  \href{https://www.caces.gob.ec/wp-content/uploads/2024/04/Modelo-generico-para-la-evaluacio\%CC\%81n-del-entorno-de-aprendizaje-de-carreras-de-grado.pdf}{PDF
  publicado por CACES}.
\item
  Localizadores: sección II, Introducción; sección V, Estructura del
  modelo y definiciones. Lectura acotada, no transcripción exhaustiva de
  indicadores.
\item
  Hallazgo pertinente: relaciona la evaluación con estándares, fuentes y
  períodos. Se utiliza como respaldo contextual de N-AA-03: identificar
  el modelo y período elegidos para un proceso.
\item
  Límite: la consulta del PDF oficial no corrobora los fixtures ni
  establece qué versión debe usarse para cada carrera hoy. No impone por
  sí sola los campos técnicos del backend ni exige un algoritmo de
  autenticación.
\item
  Estado: \textbf{documento primario localizado y consultado;
  aplicabilidad al piloto por confirmar}.
\end{itemize}

\subsubsection{F-08 --- Instructivo UNL, Resolución
037-R-UNL-2022}\label{fuentes-f-08-instructivo-unl-resoluciuxf3n-037-r-unl-2022}

\begin{itemize}
\tightlist
\item
  Autor institucional: Universidad Nacional de Loja.
\item
  Identificación: resolución 037-R-UNL-2022, instructivo de
  aseguramiento de la calidad institucional, carreras y posgrados.
\item
  \href{https://www.unl.edu.ec/sites/default/files/gaceta/resolucion/2022-06/Resoluci\%C3\%B3n\%20Nro.\%20037-R-UNL-2022-\%20Instructivo\%20para\%20el\%20Aseguramiento\%20de\%20la\%20Calidad\%20de\%20la\%20UNL-signed.pdf}{PDF
  publicado en la gaceta de la UNL}.
\item
  Localizadores: artículos 1--2 (objeto/ámbito), 9--10 (comisión de
  carrera/programa) y 12 (responsabilidades de comités). La revisión no
  abarca todas sus disposiciones.
\item
  Hallazgo pertinente: identifica comisiones y comités con
  responsabilidades en autoevaluación. Respalda el contexto
  institucional de N-AA-03 y la necesidad de explicitar quién
  interviene.
\item
  Límite: no se convierte automáticamente a una comisión en el rol
  \path|coordinador|, ni a cada miembro en administrador de catálogo. No
  se certifican vigencia actual, designaciones o delegaciones del
  piloto.
\item
  Estado: \textbf{documento primario localizado y consultado;
  correspondencia de roles y vigencia por validar}.
\end{itemize}

\subsubsection{Fuentes aún
insuficientes}\label{fuentes-fuentes-auxfan-insuficientes}

No se aportaron entrevistas al requirente, actas de designación ni una
decisión institucional vigente para el piloto. Las ediciones específicas
de Medicina, Enfermería, Odontología y Derecho continúan fuera de la
corroboración de este AA2. La localización de F-07/F-08 mejora la
procedencia documental, pero no resuelve CL-011--017 ni modifica los
manifiestos originales.

\subsection{Fuentes técnicas}\label{fuentes-fuentes-tuxe9cnicas}

\subsubsection{F-09 --- RFC 9110, HTTP
Semantics}\label{fuentes-f-09-rfc-9110-http-semantics}

\begin{itemize}
\tightlist
\item
  Autores: R. Fielding, M. Nottingham y J. Reschke; junio de 2022.
\item
  \href{https://www.rfc-editor.org/rfc/rfc9110.html}{Documento oficial
  RFC Editor}.
\item
  Localizadores revisados: §§5 y 6.3--6.4; §§8.1--8.3 y 8.8.3;
  §§9.2--9.3; §§10.2.2, 11.6, 12.5.1 y 13.1.1; §15; §17.
\item
  Uso: fundamento de decisiones HTTP y revisión de compatibilidad; no
  fuente de necesidades institucionales ni de umbrales de
  rendimiento/sesión del proyecto.
\item
  Estado: \textbf{referencia técnica consultada}; conformidad completa
  de la aplicación no acreditada.
\end{itemize}

\subsubsection{F-13 --- WCAG 2.2}\label{fuentes-f-13-wcag-2.2}

\begin{itemize}
\tightlist
\item
  Autor institucional: W3C, Web Content Accessibility Guidelines 2.2.
\item
  \href{https://www.w3.org/TR/WCAG22/}{Recomendación W3C}.
\item
  Localizadores de inicio para verificación: teclado 2.1.1, foco visible
  2.4.7 y contraste 1.4.3; el nivel AA exige revisar todos los criterios
  aplicables, no solo esos tres.
\item
  Uso: delimitar verificaciones de AA-RNF-08; RNF-014 del proyecto
  expresa el nivel objetivo.
\item
  Estado: \textbf{referencia técnica consultada; accesibilidad no
  certificada}.
\end{itemize}

\subsubsection{F-14 --- RFC 5789, PATCH Method for
HTTP}\label{fuentes-f-14-rfc-5789-patch-method-for-http}

\begin{itemize}
\tightlist
\item
  Autores: L. Dusseault y J. Snell; marzo de 2010.
\item
  \href{https://www.rfc-editor.org/rfc/rfc5789.html}{Documento oficial
  RFC Editor}, sección 2.
\item
  Uso: aclarar que PATCH tiene especificación propia. No atribuir su
  definición a RFC 9110 ni suponer JSON Merge Patch sin contrato.
\item
  Estado: \textbf{referencia complementaria consultada}.
\end{itemize}

\subsection{Índice del contraste con
implementación}\label{fuentes-uxedndice-del-contraste-con-implementaciuxf3n}

Todas las rutas se expresan desde la raíz del proyecto.

{\def\LTcaptype{none} % do not increment counter
\begin{longtable}[]{@{}
  >{\raggedright\arraybackslash}p{(\linewidth - 4\tabcolsep) * \real{0.3333}}
  >{\raggedright\arraybackslash}p{(\linewidth - 4\tabcolsep) * \real{0.3333}}
  >{\raggedright\arraybackslash}p{(\linewidth - 4\tabcolsep) * \real{0.3333}}@{}}
\toprule\noalign{}
\begin{minipage}[b]{\linewidth}\raggedright
Código
\end{minipage} & \begin{minipage}[b]{\linewidth}\raggedright
Archivo exacto
\end{minipage} & \begin{minipage}[b]{\linewidth}\raggedright
Qué permite observar
\end{minipage} \\
\midrule\noalign{}
\endhead
\bottomrule\noalign{}
\endlastfoot
C-01 &
\path|backend/src/main/java/ec/edu/unl/sec/identityaccess/api/SessionController.java|
& Login, logout, consulta de identidad, DTO y códigos. \\
C-02 &
\path|backend/src/main/java/ec/edu/unl/sec/identityaccess/application/SessionService.java|
& Vigencia/revocación, hash y control de intentos. \\
C-03 &
\path|backend/src/main/java/ec/edu/unl/sec/selfassessment/api/CareerController.java|
& Alta/consulta/cambios de unidades y carreras. \\
C-04 &
\path|backend/src/main/java/ec/edu/unl/sec/selfassessment/application/CareerAdministrationService.java|
& Código único, revisión, vigencia, consulta y permisos. \\
C-05 &
\path|backend/src/main/java/ec/edu/unl/sec/selfassessment/application/SelfAssessmentService.java|
& Creación en preparación, selección publicada y acceso al proceso. \\
C-06 &
\path|backend/src/main/java/ec/edu/unl/sec/shared/api/EntityTagSupport.java|
& Formato local de ETag y precondición. \\
C-07 &
\path|backend/src/main/java/ec/edu/unl/sec/identityaccess/infrastructure/AccessProblemWriter.java|
& Construcción de errores de acceso; no añade WWW-Authenticate. \\
C-08 &
\path|backend/src/main/java/ec/edu/unl/sec/identityaccess/api/IdentityAccessExceptionHandler.java|
& Errores de autenticación; no declara desafío de autenticación. \\
C-09 &
\path|backend/src/main/java/ec/edu/unl/sec/shared/api/ApiProblemHandler.java|
& Errores comunes con correlación. \\
C-10 &
\path|backend/src/main/java/ec/edu/unl/sec/identityaccess/infrastructure/JpaIdentityAccessRepository.java|
& Conteo de fallos por ventana, sin reinicio por éxito en esa
consulta. \\
\end{longtable}
}

\subsection{Reglas de uso de las
fuentes}\label{fuentes-reglas-de-uso-de-las-fuentes}

\begin{enumerate}
\def\labelenumi{\arabic{enumi}.}
\tightlist
\item
  Un documento del equipo respalda una decisión del producto; una norma
  primaria puede respaldar contexto, pero requiere revisar
  aplicabilidad.
\item
  Una norma HTTP no justifica quién administra una carrera; un
  instructivo institucional no define automáticamente rutas o tokens.
\item
  El código muestra implementación observada; una prueba existente no
  equivale a un resultado ejecutado en esta actividad.
\item
  El estado de una fila de trazabilidad debe reflejar sus límites de
  fuente, diseño y ejecución. En esta entrega no se asigna ``Validado
  institucionalmente''.
\end{enumerate}
\section{Nota de cambios y evidencia
individual}\label{nota-nota-de-cambios-y-evidencia-individual}

{\def\LTcaptype{none} % do not increment counter
\begin{longtable}[]{@{}
  >{\raggedright\arraybackslash}p{(\linewidth - 2\tabcolsep) * \real{0.2500}}
  >{\raggedright\arraybackslash}p{(\linewidth - 2\tabcolsep) * \real{0.7500}}@{}}
\toprule\noalign{}
\begin{minipage}[b]{\linewidth}\raggedright
Identificación
\end{minipage} & \begin{minipage}[b]{\linewidth}\raggedright
Registro
\end{minipage} \\
\midrule\noalign{}
\endhead
\bottomrule\noalign{}
\endlastfoot
Estudiante & {[}PENDIENTE: nombre completo{]} \\
Proyecto & SIGEAC / SEC-UNL \\
Actividad & AA2, dossier y matriz inicial; 4,5 horas previstas \\
Documento previo &
\path|documentos-academicos/APE2-semana-2/specification.md| \\
Revisión actual & \url{DOSSIER_REQUISITOS_AA2.md} y
\url{MATRIZ_TRAZABILIDAD.md} \\
Revisión personal / fecha & {[}PENDIENTE{]} \\
\end{longtable}
}

\subsection{Principales cambios efectivamente
preparados}\label{nota-principales-cambios-efectivamente-preparados}

\begin{enumerate}
\def\labelenumi{\arabic{enumi}.}
\tightlist
\item
  Se consultaron RFC 9110, dos documentos primarios CACES/UNL y
  referencias de accesibilidad/PATCH, registrando localizadores y
  límites en F-01--14.
\item
  La historia de catálogo APE2/US2 se separó en alta, lectura, edición y
  desactivación: HU-AA-02--05. Se conservaron acceso y preparación como
  HU-AA-01/06, sin ampliar alcance.
\item
  Los diez FR locales de APE2 se desglosaron en catorce obligaciones
  para evitar que login, logout, consulta de identidad o cambios de
  carrera se consideren un único caso probado.
\item
  Se separaron seguridad, expiración y privacidad, y evento de auditoría
  frente a su inmutabilidad/correlación. Se retienen diez RNF con
  verificaciones VC-AA-01--10.
\item
  Se separaron los rechazos compuestos de APE2 en criterios por formato,
  unidad y fechas; el dossier contiene 27 AC y diez VC, no resultados de
  ejecución.
\item
  Se construyeron 24 relaciones completas en la matriz; el estado
  expresa trazado, revisión, medición o validación de dominio pendiente.
\item
  Se registró HTTP-AA-01: contraste pendiente de WWW-Authenticate en
  respuestas 401. No se inventó un esquema Bearer ni se modificó código.
\end{enumerate}

\subsection{Equivalencia de requisitos
funcionales}\label{nota-equivalencia-de-requisitos-funcionales}

{\def\LTcaptype{none} % do not increment counter
\begin{longtable}[]{@{}
  >{\raggedright\arraybackslash}p{(\linewidth - 4\tabcolsep) * \real{0.3333}}
  >{\raggedright\arraybackslash}p{(\linewidth - 4\tabcolsep) * \real{0.3333}}
  >{\raggedright\arraybackslash}p{(\linewidth - 4\tabcolsep) * \real{0.3333}}@{}}
\toprule\noalign{}
\begin{minipage}[b]{\linewidth}\raggedright
Origen APE2
\end{minipage} & \begin{minipage}[b]{\linewidth}\raggedright
Elementos AA2
\end{minipage} & \begin{minipage}[b]{\linewidth}\raggedright
Razón de la revisión
\end{minipage} \\
\midrule\noalign{}
\endhead
\bottomrule\noalign{}
\endlastfoot
APE-FR-01 & AA-RF-01/02/03 & Separar iniciar sesión, cerrarla y
consultar funciones. \\
APE-FR-02 & AA-RF-04 & Conservar autorización por capacidad/ámbito sin
repetirla como beneficio de todas las historias. \\
APE-FR-03 & AA-RF-05/06 & Separar alta de unidad y carrera y sus
validaciones. \\
APE-FR-04 & AA-RF-07 & Precisar catálogo, detalle e historial
visible. \\
APE-FR-05 & AA-RF-08/09 & Distinguir edición y desactivación; comparten
conservación histórica, no resultado. \\
APE-FR-06 & AA-RF-10 & Identificar versión publicada sin inferir
vigencia institucional. \\
APE-FR-07 & AA-RF-11 & Preparación con referencias exactas, sin
activación. \\
APE-FR-08 & AA-RF-12 & Consulta de proceso según permiso aplicable. \\
APE-FR-09 & AA-RF-13 & Distinguir registro obligatorio de la calidad de
inmutabilidad. \\
APE-FR-10 & AA-RF-14 & Mantener error estable, con
privacidad/correlación verificadas aparte. \\
\end{longtable}
}

\subsection{Equivalencia y disposición de requisitos no
funcionales}\label{nota-equivalencia-y-disposiciuxf3n-de-requisitos-no-funcionales}

{\def\LTcaptype{none} % do not increment counter
\begin{longtable}[]{@{}
  >{\raggedright\arraybackslash}p{(\linewidth - 4\tabcolsep) * \real{0.3333}}
  >{\raggedright\arraybackslash}p{(\linewidth - 4\tabcolsep) * \real{0.3333}}
  >{\raggedright\arraybackslash}p{(\linewidth - 4\tabcolsep) * \real{0.3333}}@{}}
\toprule\noalign{}
\begin{minipage}[b]{\linewidth}\raggedright
Origen APE2
\end{minipage} & \begin{minipage}[b]{\linewidth}\raggedright
Tratamiento en AA2
\end{minipage} & \begin{minipage}[b]{\linewidth}\raggedright
Justificación
\end{minipage} \\
\midrule\noalign{}
\endhead
\bottomrule\noalign{}
\endlastfoot
APE-NFR-01 & AA-RNF-01/02/03 & Hash, expiración y privacidad son
comprobaciones diferentes. \\
APE-NFR-02 & AA-RNF-04 & Se conserva el umbral del proyecto y se declara
falta de protocolo. \\
APE-NFR-03 & Global RNF-009, diferido del recorte & No hay canal
complementario operativo en el incremento para justificar su ensayo
local. No se elimina del producto. \\
APE-NFR-04 & AA-RNF-05 & Concurrencia definida como ausencia de pérdida
silenciosa, no como presencia de una cabecera. \\
APE-NFR-05 & AA-RNF-06/07 & Inmutabilidad y correlación requieren
evidencias diferentes. \\
APE-NFR-06 & AA-RNF-08 & Criterios aplicables y revisión humana; no
certificar AA por etiquetas. \\
APE-NFR-07 & AA-RNF-09 & Mantener reglas reproducibles y revisar
acoplamiento SQL. \\
APE-NFR-08 & AA-RNF-10 & Contraste de contrato y protocolo, incluida la
observación HTTP-AA-01. \\
APE-NFR-09 & Objetivo global SC-007 y P-AA-05 & No trasladar una campaña
de 20 personas al trabajo individual sin diseño ni aprobación. \\
APE-NFR-10 & Restricción de pruebas/demostración & Datos sintéticos se
mantienen como condición transversal; evita duplicar un atributo de
calidad de ejecución. \\
\end{longtable}
}

Estos cambios son una propuesta de depuración académica; no se presentan
como aprobación del equipo ni modifican los documentos previos.

\subsection{Evidencia y contribución que debe identificar el
estudiante}\label{nota-evidencia-y-contribuciuxf3n-que-debe-identificar-el-estudiante}

La preparación fue asistida por ChatGPT/Codex. Para que sea evidencia
individual, completar:

\begin{itemize}
\tightlist
\item
  Elementos revisados personalmente: {[}PENDIENTE: IDs de
  historias/requisitos/filas{]}.
\item
  Fuentes comprobadas personalmente y localizador: {[}PENDIENTE{]}.
\item
  Corrección adoptada y justificación propia: {[}PENDIENTE{]}.
\item
  Pregunta que permanece abierta y razón: {[}PENDIENTE{]}.
\item
  Fecha y enlace a revisión/commit/PR real, si se publica:
  {[}PENDIENTE{]}.
\end{itemize}

Texto de nota breve para completar: ``Revisé personalmente {[}IDs{]}. A
partir de {[}fuente/localizador{]} cambié {[}antes{]} por {[}después{]},
porque {[}motivo{]}. Relacioné el resultado con {[}AC/VC{]} y dejé
pendiente {[}decisión{]}. Mi evidencia verificable es {[}ruta o enlace
real{]}''.

El commit base \path|faf868ea5a8a85285290c06d3baf8fa49889edac|
identifica el código inspeccionado; \textbf{no acredita autoría del
AA2}. No se crearon commits, PR ni publicaciones. El manifiesto de esta
carpeta identifica integridad y línea base, no quién participó ni cuánto
tiempo trabajó.

\end{document}
